\chapter{Implementazione}
\label{chap:implementazione}

Questo capitolo descrive le scelte implementative principali del sistema, seguendo un percorso che va dagli strati più bassi verso quelli più alti, in modo da costruire progressivamente il contesto necessario alla comprensione di ciascuna parte.

Si parte dal binding tra Kotlin e la libreria nativa FMILib, illustrando come il meccanismo di interoperabilità con il codice C sia stato configurato e come le opzioni del linker siano state adattate per ciascuna piattaforma target. Su questa base si descrive poi la gestione del ciclo di vita dell'FMU, ovvero come il sistema allochi il contesto FMI, effettui il parsing dell'XML descrittivo e carichi il binario nativo. Si procede quindi con l'estrazione dei metadati e la gestione delle variabili, spiegando come le informazioni del modello vengano interrogate tramite FMILib e restituite in un formato indipendente dalla libreria. La sezione sull'esecuzione della simulazione descrive invece il loop di avanzamento temporale, la risoluzione dei value reference e la raccolta dei risultati a ogni passo.

Una sezione dedicata tratta la ricompilazione degli FMU su macOS ARM64, un caso particolare richiesto dall'architettura Apple Silicon, in cui i binari distribuiti negli FMU non sono compatibili con la piattaforma e devono essere ricompilati a partire dai sorgenti C inclusi nel pacchetto. Il capitolo si chiude con una panoramica degli strumenti di processo adottati: la configurazione Gradle multipiattaforma, la pipeline di CI/CD e il workflow di build e test su Linux, macOS e Windows.

\section{Binding C/Kotlin e configurazione del linker per piattaforma}
\label{sec:impl-binding}

L'interoperabilità tra Kotlin e la libreria nativa FMILib è realizzata tramite il meccanismo di C interop di Kotlin/Native, che genera automaticamente binding Kotlin a partire dagli header C della libreria. La configurazione del binding avviene tramite un file \texttt{.def} in cui vengono dichiarati gli header da includere e le opzioni di compilazione e linking necessarie, tra cui il percorso alla directory degli header e il nome della libreria condivisa da collegare. A partire da questa dichiarazione, il compilatore produce un modulo Kotlin che espone le funzioni e i tipi C come API direttamente utilizzabili dal codice Kotlin, boilerplate manuale minimo rappresentato dalla scrittura del file \texttt{.def}.

La libreria FMILib non è disponibile come artefatto precompilato su repository pubblici, ma viene compilata a partire dai sorgenti tramite CMake. La compilazione è gestita dal modulo \texttt{fmilib}, che si occupa di invocare CMake, eseguire l'installazione degli artefatti e renderli disponibili al resto del progetto attraverso una configurazione e un plugin Gradle dedicato. Il modulo \texttt{fmu-kt} dichiara una dipendenza su questi artefatti e li referenzia nelle opzioni di linker del target nativo, includendo sia la directory degli header sia quella delle librerie condivise prodotte dalla compilazione.

Poiché il sistema deve essere compilato su tre piattaforme distinte — Linux x64, macOS ARM64 e Windows x64 — le opzioni di linker sono configurate per piattaforma, in quanto i percorsi degli artefatti e le convenzioni del linker differiscono tra i sistemi operativi. Su Windows, ad esempio, la libreria condivisa è un file \texttt{.dll} che deve essere copiato nella stessa directory dell'eseguibile prodotto; su Linux e macOS si tratta invece di file \texttt{.so} e \texttt{.dylib} referenziabili tramite \texttt{rpath}.

Su Linux si è reso necessario un ulteriore workaround legato alla toolchain interna di Kotlin/Native. Il compilatore Konan include un proprio sysroot con una versione di glibc datata (2.19, kernel 4.9), pensata per massimizzare la compatibilità tra distribuzioni. Quando la libreria da collegare è stata compilata con una versione più recente di glibc — come accade negli ambienti CI moderni — il linker di Konan non riesce a risolvere alcuni simboli di sistema presenti nella libreria condivisa di FMILib. Per ovviare a questo problema si è fatto ricorso alle opzioni \texttt{--allow-shlib-undefined} e \texttt{--unresolved-symbols=ignore-all}, che istruiscono il linker a non fallire in presenza di simboli non risolti nelle librerie condivise, delegandone la risoluzione al linker dinamico del sistema operativo a tempo di esecuzione. Questa è la soluzione consolidata nella comunità Kotlin/Native per questo tipo di incompatibilità, accettabile in quanto i simboli in questione sono garantiti essere presenti in qualsiasi distribuzione Linux moderna su cui il binario verrà eseguito.

\section{Gestione del ciclo di vita dell'FMU}
\label{sec:impl-lifecycle}

Un FMU è distribuito come archivio ZIP con estensione \texttt{.fmu}, contenente un file XML di descrizione del modello (\texttt{modelDescription.xml}), i binari precompilati per le piattaforme supportate e, opzionalmente, i sorgenti C del modello. Prima che qualsiasi operazione possa essere eseguita sul modello, l'archivio deve essere estratto in una directory di lavoro: questa operazione è delegata direttamente a FMILib, a cui vengono forniti il percorso del file \texttt{.fmu} e la directory di destinazione dell'estrazione. La libreria si occupa autonomamente di estrarre il contenuto e di rendere disponibili i file necessari ai passi successivi.

Una volta estratto il contenuto, FMILib costruisce un contesto \ac{fmi} — una struttura opaca che rappresenta lo stato interno della sessione nella libreria — e procede all'analisi del file \texttt{modelDescription.xml}. Da questo file vengono estratti tutti i metadati del modello: nome, autore, versione, tipo di \ac{fmu} (Model Exchange, Co-Simulation o entrambi), i parametri predefiniti dell'esperimento e l'elenco delle variabili disponibili. Il risultato del parsing è una struttura dati interna a FMILib, referenziata tramite un puntatore opaco di tipo \texttt{fmi2\_import\_t}, che viene mantenuta per tutta la durata della sessione e liberata esplicitamente alla chiusura.

Il passo successivo è il caricamento del binario nativo del modello: FMILib individua la libreria condivisa appropriata per la piattaforma corrente all'interno dell'archivio estratto e la carica dinamicamente tramite le API del sistema operativo. Questo passaggio può fallire se il binario non è presente o non è compatibile con la piattaforma — caso che si presenta, ad esempio, su macOS ARM64 con FMU che distribuiscono solo binari x86\_64, e che viene gestito tramite la ricompilazione descritta nella Sezione~\ref{sec:impl-macos}.

Il ciclo di vita del modello si chiude con una fase di terminazione altrettanto esplicita: alla chiusura, il binario viene scaricato, la struttura FMI liberata e il contesto deallocato. Questa sequenza è incapsulata in \texttt{FmuLifecycleManager}, che implementa \texttt{AutoCloseable} per garantire che le risorse native vengano sempre rilasciate correttamente, anche in presenza di eccezioni.

Per quanto riguarda il ciclo di vita della simulazione, esso si innesta su quello del modello e segue le fasi definite dallo standard FMI 2.0 per la Co-Simulation: istanziamento del componente, configurazione dell'esperimento con i parametri temporali e di tolleranza, ingresso e uscita dalla modalità di inizializzazione, esecuzione dei passi di simulazione e terminazione. Queste operazioni sono gestite da \texttt{SimulationManager}, che riceve in ingresso la struttura \texttt{fmi2\_import\_t} già inizializzata da \texttt{FmuLifecycleManager} e opera su di essa per tutta la durata dell'esperimento. L'accesso alle variabili del modello durante l'esecuzione — la loro risoluzione tramite value reference e la lettura dei valori a ogni passo — è l'argomento della sezione seguente.

\section{Estrazione dei metadati e gestione delle variabili}
\label{sec:impl-metadati}

L'estrazione dei metadati del modello avviene interamente tramite le API di FMILib, che espongono funzioni dedicate per interrogare la struttura \texttt{fmi2\_import\_t} prodotta dal parsing dell'XML. Per ciascun campo — nome del modello, autore, versione, tipo di FMU e parametri predefiniti dell'esperimento — esiste una corrispondente funzione C della libreria che restituisce il valore direttamente dalla struttura interna, senza richiedere un nuovo accesso al file XML. Questi valori vengono raccolti da \texttt{InfoManagerFmu} e assemblati in un oggetto \texttt{FmuInfo}, una data class Kotlin serializzabile che costituisce il contratto tra il modulo \texttt{fmu-kt} e il backend.

La gestione delle variabili richiede un passo aggiuntivo rispetto ai semplici campi scalari. FMILib espone la lista delle variabili del modello come una struttura iterabile, da cui è possibile estrarre per ciascuna variabile il nome e il tipo base. Per le variabili di tipo reale viene anche estratta l'unità di misura, ottenuta navigando il puntatore all'unità associata alla variabile; per le variabili di tipo diverso l'unità non è applicabile e viene registrata come assente. Il risultato è una mappa da nome di variabile a unità di misura, inclusa in \texttt{FmuInfo}, che il frontend utilizza per presentare all'utente le variabili disponibili prima di avviare una simulazione. Una volta terminate le operazioni sulla lista, la struttura viene esplicitamente liberata tramite l'apposita API di FMILib, in linea con il modello di gestione manuale della memoria della libreria.

Tra i campi di \texttt{FmuInfo} figura anche il flag \texttt{canSimulate}, che indica se il modello supporta la modalità Co-Simulation. Questo valore viene derivato dal tipo di FMU rilevato durante il parsing: solo i modelli di tipo Co-Simulation o misto (Co-Simulation e Model Exchange) sono simulabili direttamente tramite FMILib, mentre i modelli di tipo Model Exchange puro richiederebbero un solver numerico esterno non previsto dal sistema.

\section{Esecuzione della simulazione}
\label{sec:impl-simulazione}

\section{Ricompilazione degli FMU su macOS ARM64}
\label{sec:impl-macos}

\section{Strumenti di processo: Gradle, CI/CD e pipeline multipiattaforma}
\label{sec:impl-processo}